\documentclass[dvisvgm,tikz,border=5pt]{standalone}
\usepackage{amsmath,amssymb}
\usepackage{CJKutf8}
\usetikzlibrary{calc,positioning}

\definecolor{themebg}{HTML}{2e362c}
\definecolor{themefg}{HTML}{edf4e8}
\definecolor{themegray}{HTML}{9ea897}
\definecolor{themecurve}{HTML}{7eb6ff}
\definecolor{themealert}{HTML}{ff7b7b}
\definecolor{themeamber}{HTML}{f5b942}

\pgfdeclarelayer{background}
\pgfdeclarelayer{foreground}
\pgfsetlayers{background,main,foreground}

\begin{document}
\begin{CJK*}{UTF8}{gbsn}
\pagecolor{themebg}
\begin{tikzpicture}[font=\small, color=themefg, text=themefg]
\tikzset{
  n/.style={circle, draw=themefg, fill=themebg, minimum size=7mm, inner sep=0pt, line width=0.8pt},
  edge/.style={themegray, line width=0.9pt}
}

\node[anchor=west, font=\bfseries] at (0,6.3) {严格 / 满 / 完全：看结构约束，不靠名称联想};
\node[anchor=west, font=\small, text=themegray] at (0,5.78)
  {三幅图只是代表性例子；这些类别并不互斥，满二叉树同时满足严格与完全。};

% Panel 1: strict but not complete
\node[font=\bfseries, text=themeamber] at (3.0,5.05) {严格，但非满、非完全};
\node[n] (a1) at (3.0,4.15) {A};
\node[n] (b1) at (1.8,3.0) {B};
\node[n] (c1) at (4.2,3.0) {C};
\node[n] (d1) at (3.5,1.75) {D};
\node[n] (e1) at (4.9,1.75) {E};
\draw[edge] (a1)--(b1); \draw[edge] (a1)--(c1); \draw[edge] (c1)--(d1); \draw[edge] (c1)--(e1);
\node[font=\small, text=themegray, align=center] at (3.0,0.75)
  {非叶节点都有两个孩子\\但最后一层不从左连续};

% Panel 2: perfect
\node[font=\bfseries, text=themecurve] at (8.3,5.05) {满（perfect）};
\node[n, draw=themecurve, fill=themecurve!10!themebg] (a2) at (8.3,4.15) {A};
\node[n] (b2) at (7.1,3.0) {B};
\node[n] (c2) at (9.5,3.0) {C};
\node[n] (d2) at (6.5,1.75) {D};
\node[n] (e2) at (7.7,1.75) {E};
\node[n] (f2) at (8.9,1.75) {F};
\node[n] (g2) at (10.1,1.75) {G};
\draw[edge] (a2)--(b2); \draw[edge] (a2)--(c2);
\draw[edge] (b2)--(d2); \draw[edge] (b2)--(e2); \draw[edge] (c2)--(f2); \draw[edge] (c2)--(g2);
\node[font=\small, text=themegray, align=center] at (8.3,0.75)
  {每一层全部填满\\所以同时严格且完全};

% Panel 3: complete but not strict
\node[font=\bfseries, text=themealert] at (13.6,5.05) {完全，但非满、非严格};
\node[n] (a3) at (13.6,4.15) {A};
\node[n] (b3) at (12.4,3.0) {B};
\node[n] (c3) at (14.8,3.0) {C};
\node[n] (d3) at (11.8,1.75) {D};
\node[n] (e3) at (13.0,1.75) {E};
\node[n] (f3) at (14.2,1.75) {F};
\draw[edge] (a3)--(b3); \draw[edge] (a3)--(c3);
\draw[edge] (b3)--(d3); \draw[edge] (b3)--(e3); \draw[edge] (c3)--(f3);
\node[font=\small, text=themegray, align=center] at (13.6,0.75)
  {最后一层从左连续\\但节点 $C$ 只有一个孩子};

\draw[themegray, line width=0.5pt] (5.65,0.35) -- (5.65,5.35);
\draw[themegray, line width=0.5pt] (11.0,0.35) -- (11.0,5.35);

\end{tikzpicture}
\end{CJK*}
\end{document}
